% ==========================================================================
%  Common macros for both PDF and web versions of the blueprint.
%  All macros here must be plasTeX-compatible.
% ==========================================================================

% -- Theorem-like environments (collected by the dependency graph) ---------
% Reset within subsections.  In a section without subsections, omit the zero
% subsection index and number directly within the section.
\newtheorem{theorem}{Theorem}[subsection]
\makeatletter
\@addtoreset{theorem}{chapter}
\@addtoreset{theorem}{section}
% Only meaningful once hyperref has defined \theHtheorem (web/PDF builds);
% skip it in bare standalone compiles (e.g. tensor-network smoke renders).
\@ifundefined{theHtheorem}{}{%
  \renewcommand*{\theHtheorem}{%
    \ifnum\value{section}=0
      \theHchapter.\arabic{theorem}%
    \else
      \ifnum\value{subsection}=0
        \theHsection.\arabic{theorem}%
      \else
        \theHsubsection.\arabic{theorem}%
      \fi
    \fi}}
\makeatother
\renewcommand{\thetheorem}{%
  \ifnum\value{section}=0
    \thechapter.\arabic{theorem}%
  \else
    \ifnum\value{subsection}=0
      \thesection.\arabic{theorem}%
    \else
      \thesubsection.\arabic{theorem}%
    \fi
  \fi}
\newtheorem{proposition}[theorem]{Proposition}
\newtheorem{lemma}[theorem]{Lemma}
\newtheorem{corollary}[theorem]{Corollary}

\theoremstyle{definition}
\newtheorem{definition}[theorem]{Definition}

\theoremstyle{remark}
\newtheorem{remark}[theorem]{Remark}
\newtheorem{example}[theorem]{Example}

% -- Number fields ---------------------------------------------------------
\newcommand{\C}{\mathbb{C}}
\newcommand{\R}{\mathbb{R}}
\newcommand{\N}{\mathbb{N}}
\newcommand{\Z}{\mathbb{Z}}
\newcommand{\Fin}[1]{\{0,\ldots,#1{-}1\}}

% -- Dirac notation (manual; braket package not available in plasTeX) ------
\newcommand{\ket}[1]{|#1\rangle}
\newcommand{\bra}[1]{\langle #1|}
\newcommand{\braket}[2]{\langle #1 | #2 \rangle}
\newcommand{\ketbra}[2]{|#1\rangle\!\langle #2|}

% -- Operators -------------------------------------------------------------
\DeclareMathOperator{\tr}{tr}
\DeclareMathOperator{\Tr}{Tr}
\DeclareMathOperator{\spn}{span}
\DeclareMathOperator{\rank}{rank}
\DeclareMathOperator{\spec}{spec}
\DeclareMathOperator{\diag}{diag}
\DeclareMathOperator{\id}{id}
\DeclareMathOperator{\mult}{mult}       % block multiplicity function

% -- Project-specific operators --------------------------------------------
\DeclareMathOperator{\algspn}{alg}       % unital algebra generated by
\DeclareMathOperator{\krausrk}{kr}       % Kraus rank
\newcommand{\wordspn}[2]{S_{#1}(#2)}    % word span S_N(A)

% -- Common shorthands -----------------------------------------------------
\newcommand{\mc}[1]{\mathcal{#1}}
\newcommand{\E}{\mathcal{E}}              % quantum channel
\newcommand{\Id}{\mathbb{1}}              % identity operator
\newcommand{\MD}{M_D(\C)}                 % D×D complex matrices
\newcommand{\MN}[1]{M_{#1}(\C)}           % n×n complex matrices
\newcommand{\GL}{\mathrm{GL}}             % general linear group
\DeclareMathOperator{\SL}{SL}             % special linear group
